\section{Konteks Capstone dan Rancangan Sistem}

\subsection{Cerita Besar Proyek}

OMNI-RAG dikembangkan sebagai proyek Capstone oleh empat orang. Tujuan proyek besar tersebut adalah menyediakan chatbot asisten hukum yang dapat menelusuri peraturan perundang-undangan Indonesia, menghubungkan ketentuan yang saling berkaitan, dan mengembalikan dasar hukum yang dapat diverifikasi. Sistem lengkap terdiri atas jalur offline untuk membentuk basis pengetahuan dan jalur online untuk menggunakan indeks serta graf ketika menerima pertanyaan. Paper ini mengambil satu bagian dari rantai tersebut, yaitu pembangunan basis pengetahuan, karena kualitas komponen offline menjadi batas atas kualitas konteks yang dapat ditemukan pada tahap berikutnya.

Pembagian tanggung jawab tim tidak diperlakukan sebagai empat aplikasi yang terpisah. Setiap kontribusi bertemu pada artefak kanonik dan kontrak layanan yang sama. Tabel~\ref{tab:peran-tim} merangkum hubungan antara kontribusi penulis dan kontribusi anggota tim lain.

\begin{table*}[t]
\centering
\caption{Pemetaan kontribusi pada proyek Capstone OMNI-RAG}
\label{tab:peran-tim}
\begin{tabularx}{0.96\textwidth}{@{}p{0.22\textwidth}p{0.31\textwidth}X@{}}
\toprule
\textbf{Kelompok} & \textbf{Ruang tanggung jawab} & \textbf{Titik integrasi}\\
\midrule
Akuisisi dan ekstraksi & Sumber dokumen, parser, pemulihan struktur, dan provenance awal. & Menghasilkan \canon{CanonicalDocument} yang menjadi masukan seragam bagi komponen pembentukan basis pengetahuan.\\
Pembentukan basis pengetahuan & \chunker, \indexer, pembangun \kg{}, materialisasi proyeksi, dan pengujian strategi. & Menghasilkan \canon{CanonicalChunks}, \canon{CanonicalIndex}, dan \canon{CanonicalGraph}; bagian ini merupakan fokus paper.\\
Retrieval dan aplikasi & Pencarian sparse, dense, atau hybrid, pemulihan konteks induk, layanan percakapan, dan autentikasi. & Mengonsumsi indeks dan graf berdasarkan identitas artefak serta \profile{} yang aktif.\\
Orkestrasi dan penjelasan & Resolusi rencana pemrosesan, dispatch, atribusi sumber, dan penyajian jejak proses. & Menggunakan metadata run, provenance, dan identitas hasil retrieval untuk menjelaskan konteks yang dipilih.\\
\bottomrule
\end{tabularx}
\end{table*}

\subsection{Arsitektur Offline}

Jalur offline mengikuti alur \textit{raw document} menuju parser, dokumen kanonik, \chunker, dan dua proyeksi yang dibangun secara paralel oleh \indexer{} dan \kg{}. Satu rangkaian tersebut dibungkus sebagai \profile{} yang menentukan datasource, versi \plugin{}, konfigurasi, dan kebijakan materialisasi. Gambar~\ref{fig:pipeline-plugin} menampilkan rancangan yang dipakai pada eksperimen.

\begin{figure*}[t]
\centering
\includegraphics[width=0.96\textwidth]{plugin-pipeline-architecture.png}
\caption{Pipeline berbasis \plugin{} dan artefak kanonik pada jalur pembentukan basis pengetahuan}
\label{fig:pipeline-plugin}
\end{figure*}

\subsection{Kontrak \plugin{} dan \profile{}}

Arsitektur modular RAG memisahkan tahap transformasi sehingga strategi dapat diuji dan diganti tanpa mengubah seluruh alur \cite{gaoModularRAG2024}. Pada OMNI-RAG, setiap \plugin{} memiliki identitas, tipe komponen, versi, entry point, kapabilitas, dan skema konfigurasi. \textit{Plugin Registry} memvalidasi manifest tersebut sebelum orkestrator memuat implementasinya. Validasi ini mencegah komponen yang menghasilkan artefak tidak kompatibel dijalankan dalam pipeline yang sama.

\profile{} adalah konfigurasi yang memilih satu implementasi untuk setiap tahap dan mengikat pipeline pada datasource tertentu. Contoh \profile{} untuk manual memilih \code{generic-structure-aware} dan \code{dense-semantic}, sedangkan \profile{} hukum memilih \code{legal-structure-aware} dan dapat mengaktifkan pembangun \kg{}. Beberapa \profile{} dapat menggunakan dokumen yang sama melalui fan-out, tetapi hasil indeks dan graf tetap diisolasi berdasarkan identitas \profile{}.

Untuk mendukung reproducibility, setiap eksekusi menyimpan identitas datasource, \profile{}, versi \plugin{}, hash konfigurasi, waktu mulai dan selesai, jumlah masukan, jumlah keluaran, jumlah penolakan, status, serta pesan kesalahan. Satu pipeline run dapat memiliki beberapa attempt ketika materialisasi gagal atau konfigurasi perlu diulang. Pencatatan tersebut menjadi lapisan monitoring operasional: peneliti dapat membandingkan kondisi eksperimen tanpa mencampurkan keluaran dari run berbeda dan operator dapat menemukan tahap yang gagal tanpa membentuk ulang semua artefak.

\subsection{Artefak Kanonik}

Kontrak data adalah inti integrasi. \canon{CanonicalDocument} mempertahankan teks, blok berhierarki, lokasi halaman, metadata, dan provenance dari parser. \chunker{} mengubahnya menjadi \canon{CanonicalChunks} yang memiliki ID stabil, jenis unit, teks, parent, posisi, dan fingerprint. Artefak yang sama menjadi masukan bagi \indexer{} dan \kg{} sehingga keduanya merujuk sumber yang identik. Tabel~\ref{tab:kontrak-artefak-paper} merangkum field dan konsumen utama.

\begin{table}[t]
\centering
\caption{Kontrak artefak pada pembentukan basis pengetahuan}
\label{tab:kontrak-artefak-paper}
\scriptsize
\begin{tabularx}{\columnwidth}{@{}l>{\raggedright\arraybackslash}X@{}}
\toprule
\textbf{Artefak} & \textbf{Isi minimal dan konsumen}\\
\midrule
\canon{CanonicalDocument} & ID dokumen, blok berhierarki, teks, halaman, metadata, dan provenance; dikonsumsi \chunker.\\
\canon{CanonicalChunks} & ID chunk, tipe retrieval/context, teks, parent, posisi, metadata, dan fingerprint; dikonsumsi \indexer{} dan \kg{}.\\
\canon{CanonicalIndex} & ID chunk, field leksikal atau vektor, filter metadata, versi model, dan provenance; dimaterialisasi ke OpenSearch.\\
\canon{CanonicalGraph} & ID simpul dan sisi, endpoint, tipe relasi, evidence, scope, dan provenance; dimaterialisasi ke Neo4j.\\
\bottomrule
\end{tabularx}
\end{table}

\subsection{Strategi \chunker{}}

\chunker{} generik menyediakan tiga strategi dengan kontrak keluaran yang sama. \textit{Sliding-window} menggunakan jendela berukuran tetap dan overlap. \textit{Recursive} mencoba pemisah dari paragraf, kalimat, kata, hingga karakter ketika batas ukuran terlampaui. \textit{Structure-aware} menggunakan outline parser, mempertahankan heading induk, dan baru melakukan pemecahan rekursif ketika bagian target terlalu panjang. Segmentasi berbasis struktur didukung oleh temuan bahwa batas semantik dapat mengurangi pemisahan informasi yang diperlukan bersama \cite{jainAutoChunker2025}.

Untuk peraturan, \textit{legal structure-aware} memusatkan unit pada Pasal dan membawa konteks Bab atau Bagian. Satu Pasal dapat menghasilkan \textit{parent chunk} yang utuh dan beberapa \textit{child chunk} yang lebih kecil. Child digunakan sebagai kandidat pencarian, sedangkan parent dipulihkan ketika konteks lengkap dibutuhkan. Pendekatan ini sesuai dengan struktur normatif yang menempatkan Pasal, Ayat, Huruf, dan Angka dalam hubungan hierarkis \cite{uuPembentukanPeraturan2011,akarajaradwongNitiBenchComprehensiveStudy2025}.

\subsection{Strategi \indexer{}}

\indexer{} hanya memproses unit bertipe retrieval atau both, tetapi tetap menyimpan unit context untuk pemulihan parent dan sitasi. \textit{Sparse indexer} menggunakan BM25 untuk mencocokkan istilah, sedangkan \textit{dense indexer} memetakan kueri dan chunk ke ruang vektor dengan BAAI/BGE-M3. \textit{Hybrid indexer} menggabungkan peringkat sparse dan dense dengan Reciprocal Rank Fusion. Perbedaan ketiga strategi ini diuji pada chunk dan datasource yang sama sehingga perubahan hasil dapat dikaitkan dengan cara pencarian, bukan dengan perubahan unit sumber \cite{robertsonProbabilisticRelevanceFramework2009,karpukhinDensePassageRetrieval2020,chenM3Embedding2024}.

\subsection{Pembangunan \kg{}}

\kg{} bersifat opsional dan menerima \canon{CanonicalChunks} yang sama. Pada domain hukum, graf deterministik membentuk tipe simpul \code{Regulation}, \code{RevisePasal}, dan \code{Pasal}. Relasi internal \code{HAS} menjaga struktur dokumen, sedangkan relasi antardokumen mencakup \code{AMENDS}, \code{ESTABLISHES}, dan \code{REVOKES}. Setiap sisi menyimpan evidence dan provenance sehingga traversal tidak terlepas dari sumbernya.

Pada tahap online, hasil retrieval langsung menjadi seed. Resolver mengelompokkan seed berdasarkan regulasi, menelusuri relasi hingga kedalaman dua, memetakan tetangga kembali ke ID chunk, lalu melakukan deduplikasi. Graf tidak menggantikan pencarian teks; graf memperluas lingkungan dokumen setelah kandidat awal ditemukan. Prinsip tersebut membedakan metrik recall langsung dari metrik jangkauan lingkungan graf, sebagaimana juga menjadi perhatian pada pendekatan GraphRAG \cite{hanRetrievalAugmentedGenerationGraphs2025}.
